% =============================================================================
% Chapter 4 --- RTL Module Design
% =============================================================================

\chapter{RTL Module Design}
\label{ch:rtl}

This chapter describes each module as an architectural unit: what it must
do, how its datapath and control logic are organized to do it, and the
trade-offs behind each major decision. The emphasis is on the
\emph{hardware block flow} (e.g.\ \texttt{PROG}~$\rightarrow$
\texttt{VERIFY}~$\rightarrow$ \texttt{ERASE}~$\rightarrow$ \texttt{RE-PROG})
that implements the closed-loop recovery policy. Cycle-accurate transition
equations, complete state graphs, and port lists are documented in the
project repository (see Supplementary Material).

\section{Parallel Interface}
\label{sec:rtl-pi}

\begin{figure}[h]
    \centering
    \includegraphics[width=0.7\textwidth]{pi_io_block.png}
    \caption{Parallel Interface I/O block diagram. The module exposes a
        host-side port set (\signame{host\_data}, \signame{host\_cmd})
        and a controller-side port set (\signame{address},
        \signame{data}, \signame{cmd}, \signame{valid}); all signals are
        combinational, so the entire block can be timed as a single
        decode stage.}
    \label{fig:pi-io-block}
\end{figure}

The parallel interface is a purely combinational host decoder. It slices
the host packet into a payload byte, a forwarded command token, a packed
binary address, and a single qualification bit \emph{valid}. \emph{valid}
asserts only when the host command is non-idle \emph{and} each of the four
one-hot fields (\textit{PE}, \textit{SA}, \textit{col\_ID},
\textit{row\_ID}) is \emph{strictly} one-hot --- all-zero and multi-hot
patterns are both rejected. \Cref{fig:host-address-split} shows how the
32-bit host transaction word is sliced at the interface boundary, and
\Cref{fig:translated-address} shows the packed binary address that the
downstream controller actually sees.

\begin{figure}[h]
    \centering
    \includegraphics[width=0.8\textwidth]{host_address_split.png}
    \caption{32-bit host transaction word split. The high byte carries
        the payload (weight nibble or pixel byte); the remaining
        24~bits are partitioned into four concatenated one-hot fields
        --- \textit{PE} (4~bits), \textit{SA} (4~bits), \textit{col\_ID}
        (8~bits), and \textit{row\_ID} (8~bits) --- each of which must
        be \emph{strictly} one-hot for the interface to assert
        \signame{valid}.}
    \label{fig:host-address-split}
\end{figure}

\begin{figure}[h]
    \centering
    \includegraphics[width=0.8\textwidth]{translated_address.png}
    \caption{Translated, packed address format passed to the controller.
        Each one-hot field of the host word is reduced to its binary
        index, and the four indices are concatenated into a 16-bit
        packed address: \textit{PE} occupies bits $[15{:}12]$,
        \textit{SA} bits $[11{:}8]$, \textit{col\_ID} bits $[7{:}4]$,
        and \textit{row\_ID} bits $[3{:}0]$. Storage and arithmetic
        inside the controller use this packed form; the one-hot
        encoding is rebuilt by single-bit left-shifts only at the
        core-facing output.}
    \label{fig:translated-address}
\end{figure}

Concentrating packet validation in a single combinational stage is a
deliberate decision: it lifts the entire ``what if the address is illegal''
class of edge cases out of the downstream FSM, leaving the controller free
to assume single-cell selection whenever it begins a transaction. The tail
decoder reduces each one-hot field to its binary index in a single
combinational pass, producing a compact 16-bit packed address that the
controller stores and operates on internally. The one-hot form is rebuilt
by a single-bit left-shift per field at the controller output, only at the
very last step before the core.

\paragraph{Why translate to packed binary.}
The translation pays for itself three ways. \emph{Storage}: a binary index
needs 2--3 bits where the one-hot field needs 4--8, so the captured
address register and every downstream counter shrinks accordingly.
\emph{Arithmetic}: verify-side comparators, buffer addressing, and
internal counters treat the cell index as an integer, which is a few bits
of slicing on the packed form versus non-trivial encoders interspersed
through the FSM on the one-hot form. \emph{Validation as a one-time
event}: once the boundary has accepted the packet, the controller can
assume structural legality and does not re-check.

\section{ANN Controller}
\label{sec:rtl-controller}

\begin{figure}[h]
    \centering
    \includegraphics[width=0.7\textwidth]{controller_io_block.png}
    \caption{ANN Controller I/O block diagram. Host-side inputs
        (\signame{address}, \signame{data}, \signame{cmd},
        \signame{valid}) arrive pre-validated from the parallel
        interface; analog-side inputs (\signame{op\_done},
        \signame{weight\_read\_data}) close the program/erase and
        verify loops. The controller emits the packed core word
        (\signame{addr\_out}, \signame{data\_out}), the three-bit
        pulse-mode bus (\signame{pulses}), the buffer-control port
        set, and a small bank of status flags
        (\signame{erase\_retry\_exhausted},
        \signame{verify\_failed}).}
    \label{fig:controller-io-block}
\end{figure}

As mentioned earlier (\Cref{sec:arch-topology}), the controller is the
only sequential block in the stack: the parallel interface decodes
combinationally and the input buffer is a passive register-array store,
so every multi-cycle decision in the system --- command sequencing,
closed-loop recovery, pulse-train shaping --- is concentrated in this
module. It is built around a central combinational block that drives
all module outputs and produces every next-state value, plus two
clocked blocks that register the main FSM state, the sub-FSM states,
the pulse counters, and the captured host bus. The ``one big
combinational block, two clocked blocks'' style gives the synthesis
tool a single coherent picture of next-state and output logic and keeps
the timing of every output predictable from the same source.

\subsection{Registered Capture of Host Data}
\label{sec:rtl-controller-capture}

The host-supplied address and payload byte are latched on the transition
out of idle. Outside that single transition, both registers retain their
values, so every multi-cycle state (program, verify, read, erase) operates
on a stable snapshot of the transaction and is immune to subsequent host
activity. The inference path is the only deliberate exception: during
\texttt{COLLECT}, the live host data byte is forwarded so that the host
can stream a sequence of pixel bytes across consecutive beats.

\subsection{Address Round-Trip and Payload Multiplexing}
\label{sec:rtl-controller-addr}

When the controller is active it decomposes the captured packed address
back into the four binary indices and rebuilds the outbound transaction
word combinationally: a single-bit left-shift by each index produces the
four one-hot fields the analog core expects. In idle the outbound word is
forced to zero so that the core sees a quiet bus between operations. The
payload byte multiplexer selects the registered byte during
program/verify/read/erase (so the same value is presented for the entire
multi-cycle operation), the live host byte during the collect phase of
inference, and zero in every other state.

\subsection{Main FSM: Hardware Block Flow}
\label{sec:rtl-controller-mainfsm}

The main FSM has nine states. The combinational block applies a uniform
set of default outputs and then overrides them per-state in a
\sv{unique case} body, so that every state has a single, well-defined
set of bus values. The architectural flow is presented in two
complementary views: \Cref{fig:top-level-fsm} isolates the command-decode
decision taken in \texttt{S\_IDLE} (which selects the operation branch),
and \Cref{fig:main-fsm-body} shows the full set of nine states together
with the closed-loop recovery policy that links them. Cycle-by-cycle
transition rules are catalogued in the repository.

\begin{figure}[h]
    \centering
    \includegraphics[width=0.7\textwidth]{top_level_fsm.png}
    \caption{Command decode and state routing from \texttt{S\_IDLE}.
        While \signame{valid} is low or \signame{cmd} equals
        \texttt{CMD\_HIZ}, the FSM holds in \texttt{S\_IDLE}; otherwise
        a \sv{unique case} on \signame{cmd} routes into exactly one
        operation branch (\texttt{S\_PROGRAM}, \texttt{S\_ERASE},
        \texttt{S\_READ}, or \texttt{S\_COLLECT\_DATA}).}
    \label{fig:top-level-fsm}
\end{figure}

\begin{table}[h]
    \centering
    \caption{Top-level command decode from \texttt{S\_IDLE}. Each branch
        is entered combinationally on the first cycle that
        \signame{valid} is high; subsequent transitions of the
        destination state are documented in \Cref{tab:main-fsm-states}.}
    \label{tab:top-level-fsm-states}
    \begin{tabularx}{\linewidth}{l X X}
        \toprule
        \textbf{State} & \textbf{Purpose} &
            \textbf{Transition Condition \& Next State} \\
        \midrule
        \texttt{S\_IDLE} &
            Quiescent; combinationally decodes \signame{cmd} when
            \signame{valid} is asserted, latches the host address and
            payload, and routes into the matching operation branch. &
            \signame{!valid} or \texttt{CMD\_HIZ} $\rightarrow$
            \texttt{S\_IDLE}; \texttt{CMD\_PROG} $\rightarrow$
            \texttt{S\_PROGRAM}; \texttt{CMD\_ERASE} $\rightarrow$
            \texttt{S\_ERASE}; \texttt{CMD\_READ} $\rightarrow$
            \texttt{S\_READ}; \texttt{CMD\_INF} $\rightarrow$
            \texttt{S\_COLLECT\_DATA}. \\
        \texttt{S\_PROGRAM} (entry) &
            Closed-loop programming branch; runs the program sub-FSM
            and hands off to verify on completion. &
            See \Cref{tab:main-fsm-states}. \\
        \texttt{S\_ERASE} (entry) &
            Erase branch; runs the erase sub-FSM. Host/verify flag
            distinguishes host-initiated from verify-initiated entry. &
            See \Cref{tab:main-fsm-states}. \\
        \texttt{S\_READ} (entry) &
            Single-cell read branch; drives the read pulse train. &
            See \Cref{tab:main-fsm-states}. \\
        \texttt{S\_COLLECT\_DATA} (entry) &
            Inference branch; streams pixel bytes into the buffer and
            advances to compute and result. &
            See \Cref{tab:main-fsm-states}. \\
        \bottomrule
    \end{tabularx}
\end{table}

\begin{figure}[h]
    \centering
    \includegraphics[width=0.85\textwidth]{main_fsm.png}
    \caption{The 9 states and the closed-loop recovery policy. Edges
        out of \texttt{S\_VERIFY} and \texttt{S\_ERASE} summarize the
        decisions taken in the sub-FSMs --- match, under-programmed,
        over-programmed, retry-exhausted --- which implement the
        \texttt{PROG}~$\rightarrow$~\texttt{VERIFY}~$\rightarrow$
        \texttt{ERASE}~$\rightarrow$~\texttt{RE-PROG} recovery policy.}
    \label{fig:main-fsm-body}
\end{figure}

\begin{table}[h]
    \centering
    \caption{Main-FSM states. The
        \texttt{PROG}~$\rightarrow$~\texttt{VERIFY}~$\rightarrow$~\texttt{ERASE}~$\rightarrow$~\texttt{RE-PROG}
        recovery policy is implemented by the branches out of
        \texttt{S\_VERIFY} and \texttt{S\_ERASE}; sub-FSM-internal
        sequencing is documented in
        \Cref{tab:prog-fsm-states,tab:verify-fsm-states,tab:erase-fsm-states}.}
    \label{tab:main-fsm-states}
    \begin{tabularx}{\linewidth}{l X X}
        \toprule
        \textbf{State} & \textbf{Purpose} &
            \textbf{Transition Condition \& Next State} \\
        \midrule
        \texttt{S\_IDLE} &
            Quiescent; decodes host command and latches captured
            address/payload on the transition out of idle. &
            \signame{!valid} or \texttt{CMD\_HIZ} $\rightarrow$
            \texttt{S\_IDLE}; otherwise routes per
            \Cref{tab:top-level-fsm-states}. \\
        \texttt{S\_RESET} &
            One-cycle pulse to the core-reset output; currently
            unreachable from idle and reserved as the attachment point
            for a future global-reset command
            (\Cref{ch:limitations}). &
            Unconditional $\rightarrow$ \texttt{S\_IDLE}. \\
        \texttt{S\_PROGRAM} &
            Drives the address onto the core word and runs the program
            sub-FSM through its physical phases (\texttt{HIZ},
            \texttt{SELECT}, \texttt{WRITE}, \texttt{WAIT\_ACK},
            \texttt{COMPLETE}). &
            On \texttt{PROG\_COMPLETE} $\rightarrow$ \texttt{S\_VERIFY}
            (the closed-loop entry to verify). \\
        \texttt{S\_VERIFY} &
            Holds the buffer in weight-read mode and runs the verify
            sub-FSM; the \texttt{CHECK} state compares readback against
            the expected weight. &
            Match $\rightarrow$ \texttt{S\_IDLE};
            under-programmed and retries available $\rightarrow$
            \texttt{S\_PROGRAM} (nudge retry);
            under-programmed and retries exhausted $\rightarrow$
            \texttt{S\_ERASE} (verify flag set);
            over-programmed $\rightarrow$ \texttt{S\_ERASE} (verify
            flag set). \\
        \texttt{S\_ERASE} &
            Runs the erase sub-FSM and drives the erase encoding on
            the pulse-mode bus. Host/verify flag distinguishes the two
            entry contexts. &
            On \texttt{ERASE\_COMPLETE}, branches per
            \Cref{tab:erase-complete}: host-initiated $\rightarrow$
            \texttt{S\_IDLE}; verify-initiated with retries
            available $\rightarrow$ \texttt{S\_PROGRAM};
            verify-initiated with retries
            exhausted $\rightarrow$ \texttt{S\_IDLE}. \\
        \texttt{S\_READ} &
            Presents the captured address and the read encoding for
            the duration of the read window. &
            \signame{pulse\_done} $\rightarrow$ \texttt{S\_IDLE}. \\
        \texttt{S\_COLLECT\_DATA} &
            Streams up to eight live pixel bytes into the buffer at
            successive slots derived from the pixel counter. &
            After eight beats (pixel-count threshold) $\rightarrow$
            \texttt{S\_COMPUTE}. \\
        \texttt{S\_COMPUTE} &
            Switches the buffer to bit-serial output mode and advances
            \signame{bit\_sel} from zero to seven so the analog core
            sees one bit of each of the eight pixels in parallel per
            cycle. &
            \signame{pulse\_done} $\rightarrow$ \texttt{S\_RESULT}. \\
        \texttt{S\_RESULT} &
            One-cycle wrap-up; switches the buffer to result-out
            mode. &
            Unconditional $\rightarrow$ \texttt{S\_IDLE}. \\
        \bottomrule
    \end{tabularx}
\end{table}

Any out-of-range state value falls through the \sv{unique case}
\sv{default} arm and returns the FSM to \texttt{S\_IDLE}, providing a
fail-safe path against a glitched state register.

\subsection{Sub-FSMs and the Closed-Loop Recovery Policy}
\label{sec:rtl-controller-subfsms}

Three sub-FSMs run inside the main FSM and together implement the
closed-loop recovery policy. Their states correspond to the physically
distinct phases of a Y-Flash operation, so each sub-state has a natural
hardware interpretation rather than an abstract control role.

\subsubsection{Program Sub-FSM}
\label{sec:rtl-controller-prog}

A five-state sequence that mirrors the physical phases of a programming
operation: \texttt{HIZ} (bus settle, one cycle high-impedance)
$\rightarrow$ \texttt{SELECT} (drive address, wait for analog
\emph{op\_done} acknowledging cell selection, simultaneously write the
expected weight nibble into the buffer for the subsequent verify)
$\rightarrow$ \texttt{WRITE} (drive the program encoding, advance the
pulse counter through the configured pulse train, exit on
\emph{pulse\_done}) $\rightarrow$ \texttt{WAIT\_ACK} (idle bus while the
in-cell transient settles, exit on second \emph{op\_done})
$\rightarrow$ \texttt{COMPLETE} (hand control to the main FSM with
next-state $=$ verify, reset the sub-FSM to \texttt{HIZ}). The pulse
train of the \texttt{WRITE} state is selected by the LUT-based shaper on
the first attempt and falls back to a short single-cycle nudge train on
retries (\Cref{sec:rtl-controller-pulses}).

\begin{figure}[h]
    \centering
    \includegraphics[width=0.65\textwidth]{prog_fsm.png}
    \caption{Program sub-FSM. Five sub-states map onto the physical
        phases of a Y-Flash programming operation: \texttt{HIZ},
        \texttt{SELECT}, \texttt{WRITE}, \texttt{WAIT\_ACK},
        \texttt{COMPLETE}.}
    \label{fig:prog-subfsm-body}
\end{figure}

\begin{table}[h]
    \centering
    \caption{Program sub-FSM states. The pulse train of
        \texttt{PROG\_WRITE} is selected by the LUT-based shaper on the
        first attempt and falls back to a short single-cycle nudge
        train on retries (\Cref{sec:rtl-controller-pulses}).}
    \label{tab:prog-fsm-states}
    \begin{tabularx}{\linewidth}{l X X}
        \toprule
        \textbf{State} & \textbf{Purpose} &
            \textbf{Transition Condition \& Next State} \\
        \midrule
        \texttt{PROG\_HIZ} &
            One-cycle bus settle; pulse-mode bus held at high
            impedance. &
            Unconditional $\rightarrow$ \texttt{PROG\_SELECT}. \\
        \texttt{PROG\_SELECT} &
            Drive the address onto the core word, wait for the analog
            \signame{op\_done} acknowledging cell selection, and
            simultaneously write the expected weight nibble into the
            buffer for the subsequent verify. &
            \signame{op\_done} $\rightarrow$ \texttt{PROG\_WRITE}. \\
        \texttt{PROG\_WRITE} &
            Drive the program pulse-mode encoding while the pulse
            counter advances through the configured pulse train (LUT
            macro-repeat on the first attempt, single-cycle nudge on
            retries). &
            \signame{pulse\_done} $\rightarrow$
            \texttt{PROG\_WAIT\_ACK}. \\
        \texttt{PROG\_WAIT\_ACK} &
            Idle bus while the in-cell programming transient settles
            on the analog side. &
            Second \signame{op\_done} $\rightarrow$
            \texttt{PROG\_COMPLETE}. \\
        \texttt{PROG\_COMPLETE} &
            Hand control back to the main FSM with next main
            state $=$ \texttt{S\_VERIFY} and reset the program
            sub-FSM. &
            Unconditional $\rightarrow$ \texttt{PROG\_HIZ}
            (and main FSM $\rightarrow$ \texttt{S\_VERIFY}). \\
        \bottomrule
    \end{tabularx}
\end{table}

\subsubsection{Verify Sub-FSM and the Recovery Branch}
\label{sec:rtl-controller-verify-check}

A four-state sequence: \texttt{IDLE} (entry, with a short-circuit path
to \texttt{CHECK} for the degenerate single-cycle read window)
$\rightarrow$ \texttt{WAIT} (pulse-train active, increment index counter
each cycle until the configured length is reached) $\rightarrow$
\texttt{CHECK} (apply the recovery decision, below) $\rightarrow$
\texttt{DONE} (on success, return the main FSM to idle).

\texttt{CHECK} is the architectural heart of the controller: it compares
the analog-side readback against the expected weight in the input buffer
and branches into one of three outcomes --- match, under-programmed,
over-programmed --- that drive the entire recovery policy. The branch
into \texttt{ERASE} on an over-programmed readback follows directly from
the device physics: there is no monotonic way to subtract injected
charge with a program pulse on a Y-Flash cell, so the only valid
recovery from over-programming is to erase the cell and reprogram it
from a known low-conductance state. The detailed branch logic is
documented row-by-row in \Cref{tab:verify-fsm-states}.

The per-cell retry counter is cleared on a successful verify, on an
over-programmed readback (the failure mode is not retry exhaustion), and
on completion of a verify-initiated erase. Together these resets scope
the counter to a single cell at a time.

\begin{figure}[h]
    \centering
    \includegraphics[width=0.65\textwidth]{verify_fsm.png}
    \caption{\texttt{CHECK} state branching logic: match,
        under-programmed, over-programmed. The three outcomes drive the
        closed-loop recovery policy --- success returns to idle,
        under-programmed retries (or falls back to erase when the
        per-cell retry budget is exhausted), and over-programmed
        triggers an erase followed by reprogramming.}
    \label{fig:verify-subfsm-body}
\end{figure}

\begin{table}[h]
    \centering
    \caption{Verify sub-FSM states. The four states correspond to the
        readback pulse-train window (\texttt{IDLE}, \texttt{WAIT}), the
        comparison decision (\texttt{CHECK}), and the closing
        handshake to the main FSM (\texttt{DONE}). Branches in
        \texttt{CHECK} reach back into \texttt{S\_PROGRAM} or
        \texttt{S\_ERASE} of the main FSM
        (\Cref{tab:main-fsm-states}).}
    \label{tab:verify-fsm-states}
    \begin{tabularx}{\linewidth}{l X X}
        \toprule
        \textbf{State} & \textbf{Purpose} &
            \textbf{Transition Condition \& Next State} \\
        \midrule
        \texttt{VERIFY\_IDLE} &
            Entry state for the verify pulse-train window; provides a
            short-circuit path into \texttt{CHECK} for the degenerate
            single-cycle read configuration. &
            Pulse-train length $\le 1$ $\rightarrow$
            \texttt{VERIFY\_CHECK}; otherwise $\rightarrow$
            \texttt{VERIFY\_WAIT}. \\
        \texttt{VERIFY\_WAIT} &
            Pulse-train active; the verify index counter is incremented
            each cycle until the configured pulse-train length is
            reached. &
            Counter reaches configured length $\rightarrow$
            \texttt{VERIFY\_CHECK}. \\
        \texttt{VERIFY\_CHECK} &
            Compare the analog readback against the expected weight
            held in the buffer and branch into one of three outcomes
            that drive the recovery policy. &
            Match (readback $=$ expected) $\rightarrow$
            \texttt{VERIFY\_DONE};
            under-programmed (readback $<$ expected) with retries
            available $\rightarrow$ main FSM \texttt{S\_PROGRAM} (nudge
            retry);
            under-programmed and retries
            exhausted $\rightarrow$ main FSM \texttt{S\_ERASE} (verify
            flag set);
            over-programmed (readback $>$ expected) $\rightarrow$ main
            FSM \texttt{S\_ERASE} (verify flag set). \\
        \texttt{VERIFY\_DONE} &
            Successful verify; return the main FSM to idle and clear
            the per-cell retry counter. &
            Unconditional $\rightarrow$ \texttt{VERIFY\_IDLE} (and main
            FSM $\rightarrow$ \texttt{S\_IDLE}). \\
        \bottomrule
    \end{tabularx}
\end{table}

\subsubsection{Erase Sub-FSM}
\label{sec:rtl-controller-erase}

A five-state sequence whose structure mirrors the program sub-FSM,
reflecting the symmetric structure of the two physical operations on the
device: \texttt{HIZ} $\rightarrow$ \texttt{SELECT} $\rightarrow$
\texttt{PULSE} (erase encoding on the pulse-mode bus) $\rightarrow$
\texttt{WAIT\_ACK} (in-cell discharge transient) $\rightarrow$
\texttt{COMPLETE}. The exit from \texttt{COMPLETE} is governed by the
host/verify flag and the global retry counter; the per-state breakdown
is given in \Cref{tab:erase-fsm-states}, and the exit-decision matrix
is summarized in \Cref{tab:erase-complete}.

\begin{figure}[h]
    \centering
    \includegraphics[width=0.65\textwidth]{erase_fsm.png}
    \caption{Erase sub-FSM. The state graph is symmetric with the
        program sub-FSM (\Cref{fig:prog-subfsm-body}), reflecting the
        equivalent physical sequencing of the two device operations.
        The host/verify flag set on entry steers the exit from
        \texttt{ERASE\_COMPLETE} between a return to idle
        (host-initiated, or verify-initiated with the retry budget
        exhausted) and a re-entry into \texttt{S\_PROGRAM}
        (verify-initiated with retries available).}
    \label{fig:erase-subfsm-body}
\end{figure}

\begin{table}[h]
    \centering
    \caption{Erase sub-FSM states. The structure mirrors
        \Cref{tab:prog-fsm-states} state-for-state; the only
        substantive difference is the closing branch in
        \texttt{ERASE\_COMPLETE}, detailed in
        \Cref{tab:erase-complete}.}
    \label{tab:erase-fsm-states}
    \begin{tabularx}{\linewidth}{l X X}
        \toprule
        \textbf{State} & \textbf{Purpose} &
            \textbf{Transition Condition \& Next State} \\
        \midrule
        \texttt{ERASE\_HIZ} &
            One-cycle bus settle; pulse-mode bus held at high
            impedance. &
            Unconditional $\rightarrow$ \texttt{ERASE\_SELECT}. \\
        \texttt{ERASE\_SELECT} &
            Drive the address onto the core word and wait for the
            analog \signame{op\_done} acknowledging cell selection. &
            \signame{op\_done} $\rightarrow$ \texttt{ERASE\_PULSE}. \\
        \texttt{ERASE\_PULSE} &
            Drive the erase pulse-mode encoding while the pulse counter
            advances through the configured erase pulse train. &
            \signame{pulse\_done} $\rightarrow$
            \texttt{ERASE\_WAIT\_ACK}. \\
        \texttt{ERASE\_WAIT\_ACK} &
            Idle bus while the in-cell band-to-band tunneling discharge
            transient settles. &
            Second \signame{op\_done} $\rightarrow$
            \texttt{ERASE\_COMPLETE}. \\
        \texttt{ERASE\_COMPLETE} &
            Closing branch; selects the next main state from the
            host/verify flag and the global retry budget. &
            Host-initiated $\rightarrow$ \texttt{S\_IDLE};
            verify-initiated and retries
            available $\rightarrow$ \texttt{S\_PROGRAM};
            verify-initiated and retries
            exhausted $\rightarrow$ \texttt{S\_IDLE}. See
            \Cref{tab:erase-complete}. \\
        \bottomrule
    \end{tabularx}
\end{table}

\begin{table}[h]
    \centering
    \caption{Exit rules from \texttt{ERASE\_COMPLETE}.}
    \label{tab:erase-complete}
    \begin{tabularx}{\linewidth}{l l X}
        \toprule
        \textbf{Condition} & \textbf{Next main state} & \textbf{Effect} \\
        \midrule
        Host-initiated erase & Idle &
            One-shot host operation; no retry. \\
        Verify-initiated, retries available & Program &
            Recovery loop; re-enters the program state. Global retry
            counter increments. \\
        Verify-initiated, retries exhausted & Idle &
            Cell declared a hard failure; the erase-retry-exhausted
            flag is asserted to the upper layer. \\
        \bottomrule
    \end{tabularx}
\end{table}

\subsection{Pulse Generation and LUT-Based Programming}
\label{sec:rtl-controller-pulses}

Every active state that touches the analog core is pulse-shaped. The
controller maintains three coherent artifacts: a combinational per-state
pulse-length value, a registered cycle counter, and a combinational
pulse-completion signal. Outside the active states the pulse-length is
zero and the completion signal is held low, so no spurious completion can
ever advance the FSM.

The pulse-mode output to the core defaults to its idle (high-impedance)
encoding. For each active state, a per-state Boolean --- derived either
from a standard burst pattern or, for the LUT-based programming case,
from a macro-repeat pattern --- gates the active encoding onto the bus
during burst cycles and falls back to idle during inter-burst windows.
The bus is therefore a time-multiplexed burst train and never enters an
undefined state. The physical structure of a single burst is illustrated
in \Cref{fig:pulse-train-example}.

\begin{figure}[h]
    \centering
    \includegraphics[width=0.8\textwidth]{pulse_train_example.png}
    \caption{Physical pulse-train structure on the three-bit pulse-mode
        bus. A burst is the concatenation of \signame{pulse\_num}
        discrete active pulses, each $T_{prog}$ cycles wide, separated
        by single-cycle \texttt{HIZ} (\texttt{3'b000}) idle gaps. The
        active mode (\texttt{3'b010} for program in the illustration)
        is driven only while the gating Boolean is high; between
        pulses, and after the burst completes, the bus reverts to
        \texttt{HIZ}. Total burst length is therefore
        $\texttt{pulse\_num}\cdot T_{prog} + (\texttt{pulse\_num}-1)\cdot
        \texttt{PULSE\_GAP}$. \emph{Numeric example, LUT-driven first
        \texttt{PROG}:} for an expected weight of $w_e = 3$ the LUT
        emits a macro-repeat count $R = 6$ (see
        \texttt{weight\_pulse\_lut.mem} in the repository). With the
        baseline configuration $T_{prog} = 2$,
        $\texttt{pulse\_num} = 1$, and $\texttt{PULSE\_GAP} = 1$, the
        inner macro is two cycles long; the controller plays
        $R = 6$ back-to-back copies of that macro separated by
        single-cycle \texttt{HIZ} gaps, yielding a total
        \texttt{PROG\_WRITE} duration of
        $6\cdot 2 + 5\cdot 1 = 17$ cycles, of which $12$ are active
        (\texttt{PULSE\_MODE\_PROG} on the bus) and $5$ are idle
        \texttt{HIZ} gaps between copies.}
    \label{fig:pulse-train-example}
\end{figure}

\paragraph{Data-driven first programming attempt.}
On the first programming attempt for a given cell, the controller shapes
the program pulse based on the expected weight value. A look-up table
indexed by the four-bit expected weight returns a macro-repeat count:
how many copies of the standard program burst to concatenate, separated
by idle windows. The LUT is loaded from a memory image at elaboration
time and can be replaced without modifying the RTL --- this is the
mechanism by which the controller absorbs device-characterization data
without re-synthesis. \Cref{fig:lut-macro-repeat} shows the resulting
structure.

\begin{figure}[h]
    \centering
    \includegraphics[width=0.85\textwidth]{example-image.png}
    \caption{LUT-based macro-repeat structure for the first programming
        attempt. The inner macro is one standard burst; the controller
        concatenates a number of macros given by the LUT entry for the
        expected weight, separated by idle windows. Retry attempts on the
        same cell switch to a single-cycle nudge train, on the assumption
        that the first attempt has already moved the cell most of the
        way to target.}
    \label{fig:lut-macro-repeat}
\end{figure}

The architectural justification for the LUT path is three-fold:
per-weight pulse shaping can be re-calibrated by editing the LUT image
without revisiting the RTL, which gives the device-characterization team
a clean handoff point; the combination of a LUT-scaled first attempt
with a short retry pulse is consistent with the device-physics observation
that incremental programming pulses are more controllable than repeated
large ones~\cite{Sebastian2020IMC}; and if the LUT mechanism is disabled
at elaboration, the controller falls back to a fixed-length program
burst and the recovery policy still operates --- the first-attempt
success rate simply drops.

\subsection{Inference Path}
\label{sec:rtl-controller-inf}

The inference path is a four-state sequence: \emph{collect-data} (eight
pixel bytes streamed into the buffer at row-major addresses derived from
the running pixel count), \emph{compute} (buffer switched to bit-serial
mode, bit-select line advanced from zero to seven so the analog core sees
one bit of each of the eight pixels in parallel per cycle),
\emph{result} (one-cycle wrap-up in result-out mode), and \emph{idle}.
Unlike the program/erase paths, inference does not consume
\emph{op\_done}: it is a read-only operation with no in-cell transient
that needs an explicit acknowledgement, so the duration of the compute
phase is set entirely by the configured inference pulse length.

\begin{figure}[h]
    \centering
    \includegraphics[width=0.85\textwidth]{example-image.png}
    \caption{Inference dataflow. Eight pixel bytes are streamed in
        collect, written into the buffer in row-major order, then exposed
        bit-serially across eight parallel lanes during compute. The
        bit-select line advances from zero to seven so the analog core
        receives one bit of each pixel per cycle.}
    \label{fig:inference-dataflow}
\end{figure}

\subsection{Recovery Loop, End to End}
\label{sec:rtl-controller-recovery}

\Cref{fig:recovery-loop-conceptual} consolidates the closed-loop policy
in a single picture: programming enters verify; verify either returns to
idle on a match, re-enters programming on an under-programmed readback
within the retry budget, or enters erase on either an over-programmed
readback or retry-budget overflow; erase either returns to idle
(host-initiated, or retry-budget overflow) or re-enters programming after
a fresh erase. The diagram is the architectural answer to the
write--verify problem of \Cref{sec:bg-writeverify}: every failure mode of
an analog programming attempt is mapped to a well-defined digital
recovery path, so that the controller delivers a stable weight matrix in
the face of bounded device non-idealities.

\begin{figure}[h]
    \centering
    \includegraphics[width=0.85\textwidth]{example-image.png}
    \caption{Closed-loop recovery policy. The
        \texttt{PROG}~$\rightarrow$~\texttt{VERIFY}~$\rightarrow$
        \texttt{ERASE}~$\rightarrow$~\texttt{RE-PROG} loop maps every
        analog failure mode to a deterministic digital recovery path.}
    \label{fig:recovery-loop-conceptual}
\end{figure}

This recovery policy is the central reason that the controller is more
complex than a conventional write-path controller, and it shapes the
verification environment of \Cref{ch:verification}: every branch must be
exercised in simulation, and the only practical way to do that without a
physical analog core is to inject faults at the readback port.

\section{Input Buffer}
\label{sec:rtl-buffer}

\begin{figure}[h]
    \centering
    \includegraphics[width=0.7\textwidth]{buffer_io_block.png}
    \caption{Input Buffer I/O block diagram. The controller drives the
        mode-qualified port set (\signame{mode}, \signame{addr},
        \signame{data\_in}, \signame{bit\_sel}), and the buffer
        exposes a combinational byte-read port (\signame{data\_out}),
        eight parallel bit-serial lanes
        (\signame{bit\_serial}$[7{:}0]$), and a one-bit
        \signame{ready} qualifier. The same physical register array
        serves the inference pixel window and the verify-side
        expected-weight store.}
    \label{fig:buffer-io-block}
\end{figure}

The input buffer is a synchronous register-array store. The choice of
registers rather than an SRAM macro is deliberate at this stage of the
project: the buffer is small, the per-bit access pattern required by the
bit-serial lanes is trivial to implement on registers, and the design
remains portable across synthesis flows without depending on a memory
compiler. A production follow-on could swap the array for an SRAM macro
behind the same port set.

The buffer supports three concurrent paths qualified by a mode input:

\begin{itemize}
    \item \textbf{Write path} (data-load mode). A single-port write
          accepts one byte per cycle when the address is in range. The
          shared write address generator is reused by both the
          programming and inference paths, which keeps the controller
          $\leftrightarrow$ buffer interface narrow.
    \item \textbf{Byte read path} (compute, result-out, or weight-read
          modes). A combinational byte-read used by verify and host-read
          operations. Making this combinational keeps the verify
          comparator on a single clock edge.
    \item \textbf{Bit-serial output} (compute mode). Eight parallel
          one-bit lanes, where lane $k$ exposes bit \emph{bit\_sel} of
          the entry at address \texttt{addr}$+k$. The controller drives
          \emph{bit\_sel} from zero to seven during the compute phase,
          so the array sees the buffer's eight-pixel window LSB-first
          across eight cycles.
\end{itemize}

A one-bit \emph{ready} output summarizes access legality (address in
range \emph{and} mode matches direction), giving the controller a clean
cycle-by-cycle indication that the current transaction was accepted.
